\documentclass[11pt]{article}
\usepackage[utf8]{inputenc}
\usepackage{amsmath,amsthm,booktabs}
\newtheorem{proposition}{Proposición}
% Generado por code/analysis/lp.py: no editar a mano.
\newcommand{\bShockHzero}{-1.11}
\newcommand{\seShockHzero}{0.08}
\newcommand{\loShockHzero}{-1.27}
\newcommand{\hiShockHzero}{-0.94}
\newcommand{\pShockHzero}{0.000}
\newcommand{\nobsHzero}{118}
\newcommand{\bInterHzero}{-0.77}
\newcommand{\seInterHzero}{0.08}
\newcommand{\loInterHzero}{-0.93}
\newcommand{\hiInterHzero}{-0.62}
\newcommand{\pInterHzero}{0.000}
\newcommand{\bShockHone}{-1.55}
\newcommand{\seShockHone}{0.20}
\newcommand{\loShockHone}{-1.94}
\newcommand{\hiShockHone}{-1.16}
\newcommand{\pShockHone}{0.000}
\newcommand{\nobsHone}{117}
\newcommand{\bInterHone}{-1.11}
\newcommand{\seInterHone}{0.20}
\newcommand{\loInterHone}{-1.50}
\newcommand{\hiInterHone}{-0.71}
\newcommand{\pInterHone}{0.000}
\newcommand{\bShockHtwo}{-1.53}
\newcommand{\seShockHtwo}{0.28}
\newcommand{\loShockHtwo}{-2.08}
\newcommand{\hiShockHtwo}{-0.98}
\newcommand{\pShockHtwo}{0.000}
\newcommand{\nobsHtwo}{116}
\newcommand{\bInterHtwo}{-1.10}
\newcommand{\seInterHtwo}{0.29}
\newcommand{\loInterHtwo}{-1.68}
\newcommand{\hiInterHtwo}{-0.52}
\newcommand{\pInterHtwo}{0.000}
\newcommand{\bShockHthree}{-1.27}
\newcommand{\seShockHthree}{0.29}
\newcommand{\loShockHthree}{-1.84}
\newcommand{\hiShockHthree}{-0.71}
\newcommand{\pShockHthree}{0.000}
\newcommand{\nobsHthree}{115}
\newcommand{\bInterHthree}{-1.02}
\newcommand{\seInterHthree}{0.31}
\newcommand{\loInterHthree}{-1.63}
\newcommand{\hiInterHthree}{-0.41}
\newcommand{\pInterHthree}{0.001}
\newcommand{\bShockHfour}{-1.02}
\newcommand{\seShockHfour}{0.32}
\newcommand{\loShockHfour}{-1.64}
\newcommand{\hiShockHfour}{-0.39}
\newcommand{\pShockHfour}{0.001}
\newcommand{\nobsHfour}{114}
\newcommand{\bInterHfour}{-0.97}
\newcommand{\seInterHfour}{0.32}
\newcommand{\loInterHfour}{-1.60}
\newcommand{\hiInterHfour}{-0.34}
\newcommand{\pInterHfour}{0.001}


\title{Restricción de crédito y transmisión monetaria a la inversión}
\date{Borrador}

\begin{document}
\maketitle

\section{Pregunta}
¿Reduce un endurecimiento monetario la inversión no residencial más cuando el crédito bancario está restringido?
Estimamos proyecciones locales de la inversión ante sorpresas monetarias de alta frecuencia, con una interacción
con el estado del crédito del trimestre anterior.

\section{Teoría}
Una empresa con tecnología $A k^\alpha$ afronta un coste del capital $c = r(1 + m s) + \delta$, donde $s \geq 0$ mide
el endurecimiento del crédito y $m > 0$ la prima de financiación externa. El capital óptimo es
$k^* = (\alpha A / c)^{1/(1-\alpha)}$.

\begin{proposition}
La semielasticidad de la inversión al tipo de interés es más negativa cuanto más restringido está el crédito:
\[
  \frac{\partial^2 \log k^*}{\partial r \, \partial s} = -\frac{m \delta}{(1-\alpha)\, c^2} < 0 .
\]
\end{proposition}

\section{Resultados}
La especificación prerregistrada es la proyección local a $h = 1$ (dos trimestres) con errores de Newey-West de
dos retardos ($N = \nobsHone$). Ante una sorpresa contractiva de una desviación típica, la
inversión acumulada en dos trimestres varía $\bShockHone$ puntos logarítmicos con el crédito en su media. La
interacción es $c_1 = \bInterHone$ (e.e. $\seInterHone$; intervalo de confianza $[\loInterHone, \hiInterHone]$;
$p$ unilateral $= \pInterHone$): con el crédito una desviación típica más restringido, la caída es mayor.

\begin{table}[h]
\centering
\begin{tabular}{lccccc}
\toprule
$h$ & 0 & 1 & 2 & 3 & 4 \\
\midrule
$b_h$ (choque) & $\bShockHzero$ & $\bShockHone$ & $\bShockHtwo$ & $\bShockHthree$ & $\bShockHfour$ \\
 & ($\seShockHzero$) & ($\seShockHone$) & ($\seShockHtwo$) & ($\seShockHthree$) & ($\seShockHfour$) \\
$c_h$ (choque $\times$ crédito) & $\bInterHzero$ & $\bInterHone$ & $\bInterHtwo$ & $\bInterHthree$ & $\bInterHfour$ \\
 & ($\seInterHzero$) & ($\seInterHone$) & ($\seInterHtwo$) & ($\seInterHthree$) & ($\seInterHfour$) \\
\bottomrule
\end{tabular}
\caption{Proyecciones locales de la inversión acumulada. Errores de Newey-West entre paréntesis.}
\end{table}

Un Monte Carlo con el mismo estimador y el mismo tamaño muestral (\texttt{output/MC\_REPORT.md}) confirma que el sesgo
es pequeño y que los intervalos de Newey-West tienen la cobertura nominal, así que los usamos para la inferencia.

\section{Conclusión}
La transmisión de la política monetaria a la inversión es más fuerte cuando el crédito está restringido, como predice
la Proposición 1.

\end{document}
